\uuid{5n5F}
\titre{Produit}
\auteur{bodin ; hanani ; mzari}
\organisation{exo7}
\licence{CC BY-NC-SA 4.0}
\datecreate{2020-01-20}
\niveau{L1}
\module{Algèbre}
\chapitre{Polynôme, fraction rationnelle}
\sousChapitre{Définition, degré, produit}
\difficulte{3}
\theme{Polynômes, produit}
\exo7qcmid{lille-1.polynomes.008}

\contenu{
\begin{question}
\qclassification{105.05}
\qtags{motcle=produit}
Soit $A(X) = \sum_{i=0}^n a_i X^i$.
Soit $B(X) = \sum_{j=0}^m b_j X^j$.
Soit $C(X) = A(X) \times B(X) = \sum_{k=0}^{m+n} c_k X^k$.
Quelles sont les assertions vraies ?
\begin{answers}
    \bad{$c_k = a_k b_k$}

    \good{$c_k = \sum_{i+j=k} a_ib_j$}

    \bad{$c_k = \sum_{i=0}^k a_ib_i$}

    \good{$c_k = \sum_{i=0}^k a_ib_{k-i}$}
\end{answers}
\begin{explanations}
La formule (à connaître) est
$$c_k = \sum_{i+j=k} a_ib_j = \sum_{i=0}^k a_ib_{k-i}.$$
\end{explanations}
\end{question}
}
